% Part: proof-theory
% Chapter: normalization
% Section: normalization-thm

\documentclass[../../../include/open-logic-section]{subfiles}

\begin{document}

\olfileid[id]{pt}{nor}{thm}

\olsection{Teorema Normalisasi}

!!^a{proof} berada dalam bentuk \emph{normal} jika tidak mempunyai segmen
potongan. Jelas bahwa hal ini berlaku jika dan hanya jika kita tidak dapat
menerapkan konversi permutasi maupun konversi reduksi kepadanya.
\emph{Menormalisasi} !!a{proof} berarti mengubahnya menjadi bukti normal
dengan menerapkan konversi permutasi dan reduksi sampai tidak ada potongan
yang tersisa. Bahwa hal ini selalu mungkin merupakan isi
\emph{teorema normalisasi}.

\begin{thm}\ollabel{thm:normalization} Setiap !!{proof}~\Log{N1i}
$\delta$ dapat dinormalisasi, yakni dapat diubah menjadi
!!a{proof}~$\delta'$ tanpa potongan.
\end{thm}

\begin{proof}
  Lakukan induksi pada peringkat dan panjang potongan~$\delta$. Hipotesis
  induksinya ialah bahwa setiap !!{proof} yang mempunyai peringkat potongan
  lebih rendah, atau peringkat potongan sama tetapi panjang potongan lebih
  rendah, dapat dinormalisasi.

  Jika panjang potongannya~$0$, tidak ada potongan dan $\delta$ sudah
  normal. Jadi, andaikan peringkat dan panjang potongan~$\delta$ keduanya
  $>0$. Pilih segmen potongan maksimal yang paling atas dan paling kanan.
  Jika panjang segmen potongan itu $>1$, terapkan konversi permutasi. Jika
  panjangnya~$1$, terapkan konversi reduksi yang bersesuaian. Hasilnya adalah
  !!a{proof} baru yang mempunyai peringkat potongan sama tetapi panjang
  potongan lebih rendah, atau peringkat potongan lebih rendah; karena itu,
  hipotesis induksi berlaku.
\end{proof}

Strategi yang digunakan dalam bukti di atas (selalu menangani potongan
maksimal yang paling atas dan paling kanan) menyarankan suatu urutan khusus
bagi penerapan konversi pada !!a{proof} agar menghasilkan !!{proof} normal.
Yang cukup mengejutkan, urutannya ternyata tidak penting. Penerapan konversi
permutasi dan reduksi dalam urutan apa pun pada akhirnya menghasilkan bukti
normal.

Definisi segmen dan potongan dapat dengan mudah dirumuskan bagi \Log{N2i},
dan konversi permutasi serta reduksi bagi \Log{N1i} dapat diterjemahkan
langsung ke \Log{N2i}. Akibatnya, teorema normalisasi juga berlaku bagi
\Log{N2i}.

\begin{thm}\ollabel{thm:normalization-N2i} Setiap !!{proof}~\Log{N2i}
$\delta$ dapat dinormalisasi, yakni dapat diubah menjadi
!!a{proof}~$\delta'$ tanpa potongan.
\end{thm}

\end{document}
